\documentclass[aspectratio=169,10pt,t]{beamer}
\usepackage{ac1}
\aulainfo{05}{Bases da Programação}{Identificadores, dados numéricos, tipos e constantes}

\begin{document}
\capa

\begin{frame}{Roteiro}
  \begin{itemize}
    \item Identificadores e as regras do Python (Cap.~2.6)
    \item Dados numéricos: inteiros, grandes inteiros, ponto flutuante (2.7)
    \item Dados não numéricos: texto e lógicos
    \item Tipos de variáveis; tipagem dinâmica (2.8)
    \item Constantes (2.9)
  \end{itemize}
\end{frame}

\begin{frame}{Variáveis e identificadores}
  \begin{itemize}
    \item \textbf{Variável:} nome que se refere a um valor guardado na memória
    \item \textbf{Identificador:} o nome escolhido pelo programador (de variáveis, funções, módulos\ldots)
  \end{itemize}
  \begin{block}{Regras em Python}
    \begin{itemize}\small
      \item Começa com letra ou \texttt{\_}; depois letras, dígitos ou \texttt{\_}
      \item Sem espaços ou símbolos como \texttt{-}, \texttt{\$}, \texttt{@}
      \item Diferencia maiúsculas de minúsculas: \texttt{nota} \(\neq\) \texttt{Nota}
      \item Não pode ser uma \emph{palavra reservada}: \texttt{if}, \texttt{for}, \texttt{while}, \texttt{def}, \texttt{True}\ldots
    \end{itemize}
  \end{block}
\end{frame}

\begin{frame}{Identificadores válidos ou não?}
  \centering\small
  \begin{tabular}{lcl}
    \toprule
    Identificador & Válido? & Motivo\\
    \midrule
    \texttt{media\_final} & sim & \\
    \texttt{nota1} & sim & dígito não é o primeiro caractere\\
    \texttt{1nota} & \alert{não} & começa com dígito\\
    \texttt{area do circulo} & \alert{não} & espaços\\
    \texttt{valor-total} & \alert{não} & \texttt{-} é o operador de subtração\\
    \texttt{for} & \alert{não} & palavra reservada\\
    \texttt{Nota} e \texttt{nota} & sim & são variáveis \emph{diferentes}\\
    \bottomrule
  \end{tabular}
  \vspace{2ex}

  \footnotesize Boa prática: nomes \emph{descritivos} (\texttt{area\_triang}, e não \texttt{x1}). Exercício 2.8: reescrever identificadores incorretos.
\end{frame}

\begin{frame}[fragile]{Atribuição}
\begin{lstlisting}
base = 5            # cria a variável base
altura = 6
area_triang = (base * altura) / 2
base = base + 1     # atualiza o valor
a, b = 3, 4         # atribuição múltipla
a, b = b, a         # troca de valores
\end{lstlisting}
  \begin{itemize}
    \item \texttt{=} \alert{não é igualdade} matemática: é \emph{atribuição} (calcula o lado direito e guarda no esquerdo)
    \item \texttt{base = base + 1} é válido: lê o valor antigo, soma, guarda o novo
  \end{itemize}
  \fonte{code/Cap2/prog.2.5\_area\_triangulo\_v3.py}
\end{frame}

\begin{frame}{Dados numéricos}
  \begin{columns}[T]
    \column{0.5\textwidth}
    \begin{block}{\texttt{int}}
      Inteiros: \texttt{42}, \texttt{-7}, \texttt{0}.\\
      \textbf{Sem limite fixo} em Python!
    \end{block}
    \column{0.5\textwidth}
    \begin{block}{\texttt{float}}
      Reais (ponto flutuante): \texttt{3.14}, \texttt{-0.5}, \texttt{1e-3}, \texttt{2.5e8}.
    \end{block}
  \end{columns}
  \vspace{1ex}
  \centering\small
  \begin{tabular}{lll}
    \toprule
    Operador & Operação & Exemplo\\
    \midrule
    \texttt{+ - *} & soma, subtração, multiplicação & \texttt{3 * 4} \(\to\) \texttt{12}\\
    \texttt{/} & divisão real & \texttt{7 / 2} \(\to\) \texttt{3.5}\\
    \texttt{//} & divisão inteira & \texttt{7 // 2} \(\to\) \texttt{3}\\
    \texttt{\%} & resto & \texttt{7 \% 2} \(\to\) \texttt{1}\\
    \texttt{**} & potência & \texttt{2 ** 10} \(\to\) \texttt{1024}\\
    \bottomrule
  \end{tabular}
\end{frame}

\begin{frame}[fragile]{Grandes inteiros: o Universo em milímetros}
\begin{lstlisting}
# Tamanho do Universo observável em milímetros
tamanho_universo = 93000000000        # anos-luz
segundos_por_min = 60
segundos_por_h = 60 * segundos_por_min
segundos_por_dia = 24 * segundos_por_h
segundos_por_ano = 365 * segundos_por_dia
velocidade_luz_em_mms = 299792458000  # mm/s
tamanho_universo_mm = tamanho_universo * velocidade_luz_em_mms * segundos_por_ano
print("O Universo tem aproximadamente", tamanho_universo_mm, "milímetros.")
\end{lstlisting}
  \begin{itemize}
    \item O resultado tem 27 dígitos; em C ou Java isso daria \emph{overflow}
    \item Variantes: googol \(=10^{100}\), googolplexo, folha de papel dobrada 107 vezes (Exercícios 2.11 a 2.15)
  \end{itemize}
  \fonte{code/Cap2/prog.2.6\_universo.py}
\end{frame}

\begin{frame}[fragile]{Ponto flutuante: cuidado com a precisão}
\begin{lstlisting}
print(0.1, "+", 0.2, "=", 0.1 + 0.2)
\end{lstlisting}
\begin{lstlisting}[style=saida]
0.1 + 0.2 = 0.30000000000000004
\end{lstlisting}
  \begin{itemize}
    \item Já vimos na Aula 2: 0,1 não tem representação binária exata
    \item Para \emph{mostrar} resultados use formatação: \texttt{"\{:.2f\}".format(x)} ou \texttt{f"\{x:.2f\}"}
    \item Compare \texttt{float} com tolerância, nunca com \texttt{==}
  \end{itemize}
  \fonte{code/Cap2/prog.2.7\_erropf.py}
\end{frame}

\begin{frame}{Dados não numéricos}
  \begin{columns}[T]
    \column{0.5\textwidth}
    \begin{block}{\texttt{str} (texto)}
      Entre aspas simples ou duplas:\\
      \texttt{'Olá'}, \texttt{"Python"}\\
      \texttt{'42'} é texto, não número!
    \end{block}
    \column{0.5\textwidth}
    \begin{block}{\texttt{bool} (lógico)}
      Só dois valores: \texttt{True} e \texttt{False}.\\
      Resultado de comparações.
    \end{block}
  \end{columns}
  \vspace{1ex}
  \begin{exampleblock}{Também existe \texttt{complex}}
    Números complexos: \texttt{3 + 4j}. Útil em Engenharia Elétrica (impedâncias, fasores).
  \end{exampleblock}
\end{frame}

\begin{frame}[fragile]{Comparações e valores lógicos}
\begin{lstlisting}
x = 1; y = 2; z = 1
print("x == y?", x == y)     # False
print("x == z?", x == z)     # True
print("x != z?", x != z)     # False
print("x < y?", x < y)       # True
print("y >= z?", y >= z)     # True
x = 2
print("1 < x < 3:", 1 < x < 3)   # True: comparação encadeada
\end{lstlisting}
  \begin{itemize}
    \item Operadores: \texttt{==}, \texttt{!=}, \texttt{<}, \texttt{<=}, \texttt{>}, \texttt{>=}
    \item Atenção: \texttt{=} atribui; \texttt{==} compara!
  \end{itemize}
  \fonte{code/Cap2/prog.2.8\_compara.py, prog.2.9\_compara\_v2.py}
\end{frame}

\begin{frame}[fragile]{Tipos e tipagem dinâmica}
\begin{lstlisting}
x = 23
print(x, type(x))     # 23 <class 'int'>
x = 12.4
print(x, type(x))     # 12.4 <class 'float'>
x = "Teste"
print(x, type(x))     # Teste <class 'str'>
\end{lstlisting}
  \begin{itemize}
    \item Em Python \textbf{o tipo pertence ao valor}, não à variável: a mesma variável pode apontar para tipos diferentes
    \item Conversões explícitas: \texttt{int("42")}, \texttt{float("3.5")}, \texttt{str(7)}
    \item \texttt{type(x)} mostra o tipo; \texttt{int(3.9)} \(\to\) \texttt{3} (trunca)
  \end{itemize}
  \fonte{code/Cap2/prog.2.11\_variaveis\_dinamicas.py}
\end{frame}

\begin{frame}{Constantes}
  \begin{itemize}
    \item Valores que \textbf{não devem mudar} durante o programa (\(\pi\), velocidade da luz, taxas)
    \item Python não tem constantes ``de verdade''; por \alert{convenção} usamos nomes em MAIÚSCULAS
    \item Vantagem: mudar o valor em um único lugar e deixar o código legível
  \end{itemize}
\end{frame}

\begin{frame}[fragile]{Constantes: identidade de Euler}
\begin{lstlisting}
PI = 3.1415
E = 2.71828
I = 1j
print("0 =", E**(PI*I) + 1)
print("0 = {:.2f}".format(E**(PI*I) + 1))
\end{lstlisting}
\begin{lstlisting}[style=saida]
0 = (4.490366412035485e-09+9.476672816471575e-05j)
0 = 0.00+0.00j
\end{lstlisting}
  \begin{itemize}
    \item O resultado deveria ser exatamente 0: o erro vem da aproximação de \(\pi\) e \(e\)
    \item O módulo \texttt{math} traz \texttt{math.pi} e \texttt{math.e} mais precisos (Aula 13; Exercício 2.16)
  \end{itemize}
  \fonte{code/Cap2/prog.2.12\_euler.py, prog.2.13\_eulerformatado.py}
\end{frame}

\exercicios{2}{%
  \begin{itemize}\small
    \item Identificadores: \textbf{2.8}
    \item Programas com dados numéricos: \textbf{2.9} (área do quadrado), \textbf{2.10} (segundos desde a meia-noite), \textbf{2.11 a 2.15} (Universo, googol, folha dobrada)
    \item Constantes: \textbf{2.16 a 2.18} (valores de \(\pi\) e \(e\), hipotenusa, esfera de raio 5)
  \end{itemize}}

\resumo{
  \item Identificador: letras, dígitos e \texttt{\_}; não pode começar com dígito nem ser palavra reservada
  \item Tipos básicos: \texttt{int}, \texttt{float}, \texttt{str}, \texttt{bool}, \texttt{complex}
  \item Python tem tipagem dinâmica e inteiros de tamanho ilimitado
  \item Constantes por convenção em MAIÚSCULAS
  \item Próxima aula: estrutura básica de um programa, entrada e saída
}

\end{document}
